\documentclass[10pt,letterpaper]{article}
\usepackage[T1]{fontenc}
\usepackage[utf8]{inputenc}
\usepackage[sfdefault]{sourcesanspro}
\usepackage{inconsolata}
\usepackage[margin=0.84in,headheight=15pt]{geometry}
\usepackage{microtype}
\usepackage{xcolor}
\usepackage{booktabs,longtable,tabularx,array}
\usepackage{enumitem}
\usepackage{fancyhdr,lastpage}
\usepackage[most]{tcolorbox}
\usepackage{listings}
\usepackage{titlesec}
\usepackage{hyperref}
\usepackage{amsmath,amssymb}
\usepackage{multicol}
\definecolor{navy}{HTML}{15324C}
\definecolor{teal}{HTML}{087E83}
\definecolor{slate}{HTML}{536879}
\definecolor{paper}{HTML}{F2F6F8}
\definecolor{warm}{HTML}{E8F4F3}
\hypersetup{colorlinks=true,urlcolor=teal,linkcolor=navy,pdfauthor={Software Engineering Reference Suite},pdfsubject={MISRA C++:2023 engineering guidance}}
\titleformat{\section}{\Large\bfseries\color{navy}}{\thesection}{0.7em}{}
\titleformat{\subsection}{\large\bfseries\color{teal}}{\thesubsection}{0.6em}{}
\titleformat{\subsubsection}{\normalsize\bfseries\color{navy}}{\thesubsubsection}{0.5em}{}
\setlist[itemize]{leftmargin=*,topsep=2pt,itemsep=2pt}
\setlist[enumerate]{leftmargin=*,topsep=2pt,itemsep=3pt}
\setlength{\parskip}{5pt}
\setlength{\parindent}{0pt}
\setlength{\emergencystretch}{2em}
\pagestyle{fancy}
\fancyhf{}
\fancyhead[L]{\small\color{slate}MISRA C++:2023 Engineering Suite}
\fancyhead[R]{\small\color{slate}\nouppercase{\leftmark}}
\fancyfoot[L]{\scriptsize\color{slate}Independent engineering guidance; see licensed source}
\fancyfoot[R]{\scriptsize\color{slate}\thepage\ / \pageref{LastPage}}
\newtcolorbox{sourcebox}{colback=paper,colframe=navy!30,boxrule=0.45pt,arc=1.6pt,breakable}
\newtcolorbox{practicebox}{colback=warm,colframe=teal!40,boxrule=0.45pt,arc=1.6pt,breakable}
\newcommand{\SourceNote}[1]{\begin{sourcebox}\textbf{SOURCE BASIS / NORMATIVE POINTER}\quad #1\end{sourcebox}}
\newcommand{\PracticeNote}[1]{\begin{practicebox}\textbf{ENGINEERING RECOMMENDATION / NOT MISRA TEXT}\quad #1\end{practicebox}}
\newcommand{\Record}[1]{\textsf{\bfseries #1}}
\lstdefinestyle{cppsuite}{language=C++,basicstyle=\ttfamily\footnotesize,keywordstyle=\color{navy}\bfseries,commentstyle=\color{slate},stringstyle=\color{teal},breaklines=true,breakatwhitespace=true,columns=fullflexible,frame=single,framerule=0.3pt,rulecolor=\color{navy!35},backgroundcolor=\color{paper},showstringspaces=false,tabsize=2,xleftmargin=0.5em,xrightmargin=0.5em,aboveskip=4pt,belowskip=4pt}
\lstset{style=cppsuite}
\newcommand{\SuiteTitle}[3]{%
\begin{titlepage}\thispagestyle{empty}
\vspace*{1.5cm}
{\color{teal}\rule{2.4cm}{2pt}}\par\vspace{1cm}
{\Huge\bfseries\color{navy}MISRA C++:2023\\[0.4em]#1}\par\vspace{0.7cm}
{\Large\color{slate}#2}\par\vspace{1.2cm}
\begin{sourcebox}\textbf{Source baseline:} MISRA C++:2023, Guidelines for the use of C++17 in critical systems, October 2023 (user-supplied licensed copy). This is an independent companion, not an official MISRA publication, rule replacement, certified checker, or claim of project compliance.\end{sourcebox}
\vfill{\large\color{navy}Technical edition \textbullet\ 16 October 2026}\par\smallskip
{\small\color{slate}#3}\par\end{titlepage}
\tableofcontents\clearpage
}
\newcommand{\StandardsBoundary}{%
\section*{How to read normative status}\addcontentsline{toc}{section}{How to read normative status}
\textbf{Normative source:} the licensed MISRA C++:2023 guideline text, including any applicable \emph{Amplification} and \emph{Exception}. A guideline \emph{headline} alone is not a complete compliance test. The source's Section~3 defines classifications, categories, decidability, analysis scope, and applicability.\par
\textbf{Independent guidance:} all review questions, programming patterns, lifecycle processes, sample thresholds, ownership assignments, controls, and CI configurations in this suite are \emph{engineering recommendations}, not statements issued by MISRA.\par
\textbf{Critical limitation:} no individual example is asserted compliant with the whole standard. For formal compliance claims, follow the applicable edition of \emph{MISRA Compliance} as the source itself instructs, with a documented compliance plan and supporting evidence.\par
}

\begin{document}
\SuiteTitle{Program Orientation and Adoption}{Scope, normative boundaries, responsibility model, and introduction}{Volume 01 / Program}
\StandardsBoundary

\section{Purpose, intended users, and scope}
This program establishes a controlled C++17 subset for teams building high-integrity embedded or hosted software. The working baseline is ISO/IEC 14882:2017 as referenced by MISRA C++:2023 Section 1.3. C++20/23/26 features are outside this baseline unless individually evaluated and documented.

The companion suite provides a catalog, rule-wise learning material, coding fixtures, analysis strategy, compliance operating procedures, deviation controls, CI example architecture, training, and evidence templates. Do not use the suite alone to determine the conformance of a specific program.


\section{What the standard distinguishes}
\small\setlength{\tabcolsep}{3pt}\renewcommand{\arraystretch}{1.15}
\begin{longtable}{p{0.22\linewidth}p{0.35\linewidth}p{0.35\linewidth}}
\toprule
Dimension & Source-aligned meaning & Practical implication\\ \midrule
\endfirsthead
\toprule
Dimension & Source-aligned meaning & Practical implication\\ \midrule
\endhead
Rule vs directive & A rule states a source-checkable requirement; a directive may require facts beyond source code. & Create code analyzer coverage for rules, and design/process evidence for directives.\\
Mandatory & No violations and no deviations are permitted. & Block unresolved actual violations; correction is necessary.\\
Required & Meet the guideline or produce an authorized formal deviation. & Use a controlled exception approval, with rationale and evidence.\\
Advisory & Meet where practical or document the reason for an observed violation. & A record is still expected; advisory is not equivalent to ignore.\\
Disapplied & Possible project recategorization of an advisory guideline, with justification. & Do not silently disable diagnostics.\\
Decidable vs undecidable & Classification of the theoretical computability of an analysis. & For undecidable cases explicitly handle uncertainty and analysis limits.\\
Single Translation Unit vs System & Amount of program needed for sufficient analysis. & Whole-program analyses need complete program variants and linker contexts.\\
\bottomrule
\end{longtable}
\normalsize
\SourceNote{Source: Sections 1.3--1.5 and 3.1--3.8, notably Section 3.8 on normative amplifications and applicable exceptions. The standard refers compliance claims to the current edition of MISRA Compliance.}

\section{Adoption in a critical software program}
\begin{enumerate}
\item Designate a baseline: C++17 dialect, compiler/toolchain version, target architecture, build flags, translation-unit inventory, product variants, and third-party boundary.
\item Create a guideline applicability map: include all 179 identifiers; preserve directive/rule, category, decidability, scope, and official source locator.
\item Select a primary analyzer supporting the intended C++17 subset. Perform a rule-by-rule coverage assessment rather than assuming a vendor checkset is exhaustive.
\item Capture documentation required for directives and cross-translation-unit rules; assign a responsible engineer and independent reviewer.
\item Build with reproducible configurations; run initial scans; triage actual violations separately from tool uncertainty, known limitations and source-defined exceptions.
\item Resolve mandatory violations; remediate or formally deviate Required violations; document Advisory violations or justified disapplications.
\item Review project evidence before stating compliance; repeat after meaningful changes in code, compilation flags, target compiler or analyzer configurations.
\end{enumerate}

\section{Baseline and legacy decisions}
Section 1.4 notes that retrofitting legacy code may create additional defects. Tailor the rollout by risk and change scope rather than mass-editing trusted modules without regression tests. Any resulting compliance claim must nonetheless reflect actual remaining violations and their treatment.

Recommended staged adoption: new code first; safety-critical modified paths second; stabilized legacy code after coverage-based impact assessment. Record temporary technical debt separately from formal deviations, because debt tickets do not authorize violations.


\section{Roles and separation of duties}
\small\setlength{\tabcolsep}{3pt}\renewcommand{\arraystretch}{1.15}
\begin{longtable}{p{0.25\linewidth}p{0.63\linewidth}}
\toprule
Role & Recommended accountabilities\\ \midrule
\endfirsthead
\toprule
Role & Recommended accountabilities\\ \midrule
\endhead
System owner & Defines criticality, program scope and acceptance risk.\\
C++ lead & Owns language subset, code architecture and guideline interpretation escalations.\\
Developer & Writes safe C++17, tests and reviews analyzer findings.\\
Independent reviewer & Validates resolution and evidentiary support for findings.\\
Assurance / compliance lead & Owns guideline plan, deviations, audit package and compliance claim approval.\\
Build / DevSecOps engineer & Maintains compiler flags, analysis reproducibility and immutable pipeline evidence.\\
\bottomrule
\end{longtable}
\normalsize

\section{Required records at program kickoff}
\begin{itemize}
\item CFG-001: compiler and language options manifest with platform assumptions.
\item CMP-001: project compliance plan and scope register.
\item GDL-001: 179-line guideline applicability and classification register.
\item ANA-001: checker-to-guideline coverage and limitations matrix.
\item DEV-001: deviation and documented-advisory register.
\item EVD-001: immutable scan, review and build evidence index.
\end{itemize}

\end{document}
